\documentclass[11pt,a4paper]{article}
\usepackage[left=0.75in,right=0.75in,top=0.80in,bottom=0.85in]{geometry}
\usepackage{graphicx,amsmath,amssymb,booktabs,caption,subcaption,hyperref,xcolor,float}
\usepackage{setspace,titlesec,enumitem,multirow,tabularx, xcolor}

\setstretch{1.08}
\setlength{\parskip}{4pt}
\setlength{\parindent}{0pt}
\titlespacing*{\section}{0pt}{10pt}{5pt}
\titlespacing*{\subsection}{0pt}{7pt}{4pt}
\titlespacing*{\subsubsection}{0pt}{5pt}{3pt}
\setlist{nosep,leftmargin=*,topsep=2pt}
\captionsetup{font=small,skip=5pt}

\graphicspath{{output/}}

\title{\textbf{PCA-Based Butterfly Trading on the GBP OIS Curve:\\
Construction, Performance, and Strategy Variants}}
\author{
Christopher Darcy Morkan (06052066) \quad Ilia Bolgov (06051900) \quad Raphael Delforge\\[4pt]
\textit{Imperial College London,  Topics in Derivative Pricing}\\[2pt]
\textit{Term 2, April 2026}
}
\date{}

\begin{document}
\maketitle

% =====================================================================
\begin{abstract}
We construct and evaluate PCA-based curvature butterfly strategies on
GBP overnight indexed swap (OIS) rates from January~2023 to January~2026,
a period spanning the Bank of England's hiking, hold, and cutting
regimes.  Butterflies are formed by solving the loading-nullity system
$L^{\!\top}w = [0,0,1]$ to eliminate level and slope exposure while
targeting unit curvature sensitivity.  We compare three curve
configurations (spot 1Y--30Y, spot 1Y--10Y, and 1-year forward
rates) and find that forward-rate PCA delivers materially stronger
curvature isolation ($R^{2}=0.83$, curvature correlation $\rho=0.90$)
than spot-rate implementations ($R^{2}=0.06$--$0.17$).  A z-score
mean-reversion signal achieves a Sharpe ratio of 1.42 on the forward
curve.  We further test three literature-inspired strategy
variants (volatility-scaling, momentum, and carry overlay) and find
though no variant is universally superior across curve configurations;
results are highly regime- and parameterisation-dependent.
Results demonstrate that principal-component ordering is unstable across
regimes, motivating structural PC identification and rolling estimation.
\end{abstract}

% =====================================================================
\section{Introduction}
\label{sec:intro}
Principal component analysis (PCA) of yield-curve changes is a
foundational tool in fixed-income analytics.  Since the early work of
Litterman and Scheinkman~(1991), it has been well established that three
factors, level, slope, and curvature, explain upwards of 95\% of
term-structure variance in most developed markets.  These factors
provide a natural basis for constructing relative-value trades:  a
\emph{butterfly} trade that is hedged against level and slope movements
isolates exposure to the curvature factor, offering a systematic
risk premium for bearing higher-order term-structure risk.

However, the practical implementation of PCA butterflies raises several
non-trivial choices.  First, PCA factors are statistical artefacts whose
economic interpretation depends on the eigenvalue ordering, which in turn
depends on the covariance structure of rates in the estimation window.
As monetary-policy regimes change, and with them the relative volatility
of parallel shifts versus twist and curvature moves, the ordinal
labelling of PCs can swap.  Second, the choice of rate representation
(spot versus forward) affects factor isolation: forward rates amplify
local curvature and reduce the multicollinearity inherent in overlapping
zero-coupon rates.  Third, the estimation window, number of tenors, and
stationarity treatment all influence factor stability.

This paper addresses these issues in the context of GBP OIS rates during
2023--2026, a period that includes the Bank of England's terminal hiking
phase (Jan--Aug~2023), an extended hold at 5.25\% (Aug~2023--Aug~2024),
and the subsequent cutting cycle (Aug~2024--Jan~2026).  Our
contributions are threefold: (i)~we show that forward-rate PCA produces
materially cleaner curvature trades than spot-rate PCA; (ii)~we
demonstrate that structural PC identification (by loading shape rather
than eigenvalue rank) is essential for regime-robust trading; and
(iii)~we evaluate three literature-inspired signal and sizing variants,
finding that no strategy variant is uniformly superior across curve
configurations, and that results are concentrated in a single regime.

The remainder of the paper is organised as follows.
Section~\ref{sec:data} describes the data and curve construction.
Section~\ref{sec:pca} sets out the PCA framework and the choices made.
Section~\ref{sec:butterfly} explains the butterfly weight solver.
Section~\ref{sec:strategies} describes the trading strategies.
Section~\ref{sec:results} presents the results, including factor
correlations, regime analysis, and strategy variant comparison.
Section~\ref{sec:discussion} discusses residual exposure, regime
dependence, and the economic rationale for the observed performance.
Section~\ref{sec:conclusion} concludes.


% =====================================================================
\section{Data and Curve Construction}
\label{sec:data}

\subsection{Raw Data}
The dataset comprises 762 daily observations of GBP OIS par rates
from January~2023 to January~2026, sourced from Bloomberg. The raw
data contains 29 maturity columns at tenors ranging from overnight (is it not from 1W to 50Y) 
to 30~years. We map Bloomberg ticker suffixes to numeric maturities
using a predefined dictionary (e.g.\ \texttt{1Y}$\to$1.0,
\texttt{18M}$\to$1.5, \texttt{30Y}$\to$30.0) and convert all rates
from percentage to decimal form.

\subsection{Par Rates as Proxies for Spot Rates}
\label{sec:par_vs_zero}
Bloomberg quotes GBP OIS rates as \emph{par rates}, corresponding to
the fixed leg of an at-the-money OIS swap, rather than as bootstrapped
zero-coupon rates.  In this study, we use these par rates directly as
proxies for spot rates in both PCA estimation and butterfly construction.

This approximation is well-justified for three reasons.

\textbf{(1) The par--zero distinction is small for OIS.}
The difference between par and zero rates arises from coupon
re-investment effects.  For OIS contracts, where the floating leg
resets daily at SONIA and fixed rates are near the prevailing
overnight rate, this distortion is typically small and smooth across
maturities.  Unlike high-coupon bonds or off-market swaps, OIS par
rates closely track the underlying zero-coupon curve.

\textbf{(2) PCA and curvature strategies are robust to smooth distortions.}
Principal component analysis extracts directions of co-movement across
tenors.  A smooth maturity-dependent shift between par and zero rates
does not materially alter the relative shape of the curve, and hence
has limited impact on the identification of level, slope, and curvature
factors.  Consistent with this, both industry and academic studies
commonly apply PCA directly to quoted par swap rates without explicit
bootstrapping.

\textbf{(3) Bootstrapping introduces model dependence.}
Recovering a zero-coupon curve from par rates requires interpolation
assumptions (e.g.\ linear, spline, or parametric models), which can
distort the covariance structure of rate changes.  Using raw par rates
avoids introducing additional modelling choices beyond those already
embedded in the market data.

\textbf{Where the distinction would matter.}
For precise pricing or risk management of swap portfolios, bootstrapped
zero-coupon curves are required to compute discount factors and DV01s.
However, for a research backtest focused on relative movements and
factor structure, the par-versus-zero approximation is standard and
introduces only a second-order error relative to model and signal
uncertainty.

\subsection{Interpolation to Standard Grid}
To ensure a uniform maturity grid across all dates, we apply linear
interpolation across tenors for each cross-section to the target maturities
$\{1,2,3,5,7,10,15,20,30\}$~years.  Dates with fewer than 10 valid tenor
observations are dropped, yielding \textbf{753 clean cross-sections}.
Figure~\ref{fig:full_pca}(a) shows the full-sample PCA loadings on this grid.

\subsection{Forward Rate Computation}
One-year forward rates are computed from par rates via the
standard par-rate bootstrapping identity (exact for zero-coupon instruments
and an excellent approximation for near-at-the-money OIS):
\begin{equation}
f(T_{i-1},T_i) = \frac{T_i \, S(T_i) - T_{i-1} \, S(T_{i-1})}{T_i - T_{i-1}},
\label{eq:forward}
\end{equation}
where $S(T)$ denotes the par rate at maturity $T$ (an excellent approximation
to the true zero rate for near-at-the-money OIS; see
Section~\ref{sec:par_vs_zero}).  This yields
8 forward-rate maturities \\ $\{f(1,2), f(2,3), f(3,5), f(5,7), f(7,10), f(10,15), f(15,20), f(20,30)\}$~years, on which a
separate PCA is run.  The forward-rate representation amplifies local
curvature in the term structure, as each forward rate is a differential
rather than a cumulative measure.


\subsection{Bank of England Regime Context}
The sample spans three distinct Bank of England policy regimes:
\begin{enumerate}
  \item \textbf{Hiking} (Jan--Aug 2023): Bank Rate raised from 3.50\%
        to 5.25\%, with material daily volatility in rates.
  \item \textbf{Hold} (Aug 2023--Aug 2024): Bank Rate held at 5.25\%,
        a period of low volatility in realised short rates.
  \item \textbf{Cutting} (Aug 2024--Jan 2026): Gradual easing begins,
        with renewed volatility as expectations shift.
\end{enumerate}
These regimes create an ideal testing ground for PCA stability, as the
covariance structure of rate changes differs markedly across phases.
All plots in this paper include regime shading to facilitate visual
interpretation.


% =====================================================================
\section{PCA Framework}
\label{sec:pca}

\subsubsection{PCA on Changes, Not Levels}

We difference the curves ($\Delta y_t = y_t - y_{t-1}$) before computing
the covariance matrix.  Yield levels exhibit unit-root behaviour: PCA on
levels would be dominated by the trending component and produce spurious
``factors'' driven by the non-stationary drift.  Daily changes are
approximately stationary, making covariance-based PCA valid.  We further
\textbf{standardise} (z-score) each tenor column before PCA, so that the
sample covariance of the standardised changes equals the
\textbf{correlation matrix}. Eigen-decomposition of the correlation matrix (rather than the raw covariance matrix) reduces the likelihood that high-volatility, long-dated tenors dominate the first principal component.

\subsubsection{Estimation Window}
We use a \textbf{252-business-day rolling window} ($\approx$1~year) in all specifications. For the full Spot 1Y--30Y curve, this corresponds to a universe of $p=9$ tenors; the truncated Spot 1Y--10Y and Forward 1Y specifications use $p=6$ and $p=8$, respectively. The window is long enough to provide a stable estimation sample relative to the number of tenors, while still being short enough to adapt to regime changes. The out-of-sample period is data-dependent and is reported by the backtest at runtime.

\subsubsection{Spot vs.\ Forward Rates}

We run PCA on both spot and forward rates.  The motivation is that
forward rates isolate local term-structure curvature more effectively
than spot rates, which embed overlapping forward exposures and create
multicollinearity.

\subsubsection{Curve Segments and Tenor Choices}
Three curve configurations are tested:
\begin{itemize}
  \item \textbf{Spot 1Y--30Y} with butterfly at 3s/7s/15s. The full
        curve follows standard market convention for GBP butterflies.
  \item \textbf{Spot 1Y--10Y} with butterfly at 2s/5s/10s, a
        truncated curve that excludes the ultra-long end where
        supply--demand dynamics (pension hedging) may distort the
        curvature factor.
\item \textbf{Forward 1Y} with butterfly at 3s/7s/15s, using forward
        rates computed from the full spot curve.
\end{itemize}
The larger curve $p$ universes are used to estimate PCA factors, while the three listed tenors above define the tradable butterfly legs only. Estimating PCA factors on only three tenors would give noisy results.

\subsection{PC Identification by Shape}
A critical implementation detail is that \textbf{level, slope, and
curvature are not always PC1, PC2, and PC3}.  The standard ordering
$\lambda_1 > \lambda_2 > \lambda_3$ depends on the relative variance
of each factor in the estimation window.  During the hold regime, slope
and curvature factors can have comparable variance, causing the ordinal
assignment to flip.

We classify PCs structurally using loading-shape diagnostics:
\begin{itemize}
  \item \textbf{Level}: the PC whose loadings have the smallest
        standard deviation (most uniform across tenors).
  \item \textbf{Slope}: the PC whose loadings have the highest
        absolute correlation with maturity.
  \item \textbf{Curvature}: the remaining PC, typically exhibiting a
        U-shaped or inverted-U loading pattern.
\end{itemize}
This classification is re-evaluated at each rolling window.  We also
enforce a sign convention: level loadings sum to a positive value, slope
loadings increase with maturity, and curvature loadings have a positive
middle (belly).
Figure~\ref{fig:pc_role} shows the evolution of PC role assignments
over time.

\paragraph{Regime-Driven Factor Reordering.}
The reordering of slope and curvature factors is driven by changes in the
relative variance of different regions of the yield curve across monetary
policy regimes. During hiking and cutting phases, uncertainty about the
future path of policy rates is concentrated in the intermediate maturities
(the ``belly'' of the curve), as markets continuously reprice the timing,
speed, and terminal level of policy adjustments. In contrast, the very
short end is closely anchored to the current policy rate, while the long
end is tied to longer-term inflation expectations and term premia, and
therefore moves more gradually. This leads to disproportionately large
movements in the belly relative to both ends, generating pronounced
non-linear distortions in the curve. As a result, the curvature factor
can exhibit higher variance than the slope factor, causing it to be
ranked as PC2 rather than PC3. In more stable ("Hold") regimes, where the curve
evolves primarily through parallel shifts and tilts, slope typically
reasserts itself as the second principal component.

\subsection{Full-Sample PCA}
Before running any rolling-window or out-of-sample analysis, we first perform PCA on the entire dataset of daily spot-rate changes. This provides a reference point for the classical level--slope--curvature factor structure and confirms that our data and methodology are consistent with established literature. The high cumulative explained variance for the first three principal components justifies focusing on these factors for trading and risk analysis.

The full-sample PCA on 753 daily spot-rate changes yields explained
variance ratios of \textbf{85.1\%, 10.6\%, and 3.0\%} for the first
three components, cumulatively explaining \textbf{98.7\%} of total
variance. The loadings display the classical level--slope--curvature
shapes (Figure~\ref{fig:full_pca}). The shapes of the loadings visually confirm that the first three PCs correspond to economically meaningful curve movements:

\begin{itemize}
    \item \textbf{PC1 (Level):} The loadings are nearly flat across all maturities, indicating that this component captures parallel shifts in the yield curve. When PC1 moves, all yields tend to increase or decrease by a similar amount, which is interpreted as a change in the overall level of interest rates.
    \item \textbf{PC2 (Slope):} The loadings are negative at short maturities and positive at long maturities, increasing monotonically with maturity. This pattern means that PC2 captures changes in the steepness of the yield curve: when PC2 increases, long-term rates rise relative to short-term rates (steepening), and vice versa (flattening).
    \item \textbf{PC3 (Curvature):} The loadings are positive at the short and long ends but negative in the middle maturities, forming a “hump” or “bowl” shape. This reflects changes in the curvature of the yield curve: when PC3 increases, yields at intermediate maturities fall relative to both short and long maturities, making the curve more “humped.” Conversely, a negative PC3 means the curve is more “inverted” in the middle. This component isolates movements where the short and long ends move together in the opposite direction to the middle, which is not explained by parallel shifts or slope changes.
\end{itemize}
\begin{figure}[H]
\centering
\includegraphics[width=0.95\textwidth]{full_sample_pca.png}
\caption{Full-sample PCA results.  (a)~Factor loadings for PC1 (level),
PC2 (slope), and PC3 (curvature) on the 9-tenor spot grid.
(b)~Individual and cumulative explained variance.  The first three PCs
explain 98.7\% of total variance.}
\label{fig:full_pca}
\end{figure}

% =====================================================================
\section{Butterfly Construction}
\label{sec:butterfly}

\subsection{Methodology Overview}
A butterfly strategy is a three-tenor relative-value trade consisting of
a short maturity wing, an intermediate maturity belly, and a long
maturity wing.  In this paper, the tradable butterflies are 3s/7s/15s
for the full spot and forward curves, and 2s/5s/10s for the truncated
spot curve.  The objective is to capture changes in curve curvature,
that is, whether the belly richens or cheapens relative to the wings.

Rather than using fixed ad hoc weights, we construct the butterfly
portfolio from rolling PCA loadings.  PCA decomposes yield-curve changes
into the familiar level, slope, and curvature factors.  We then choose
weights so that the portfolio has zero exposure to level and slope and
unit exposure to curvature.  This produces a pure curvature portfolio,
whose daily PnL is driven by relative movements of the belly versus the
wings.  The resulting PCA-based butterfly is the common underlying
portfolio used in all later strategy variants; the differences across
strategies arise from how this curvature exposure is timed, scaled, or
filtered.

\subsection{Weight Solver}
Given a $3 \times 3$ matrix of PCA loadings $L$ at the three butterfly
tenors $(\tau_s, \tau_m, \tau_l)$, we solve
\begin{equation}
L^{\!\top} w = \begin{pmatrix} 0 \\ 0 \\ 1 \end{pmatrix},
\label{eq:weights}
\end{equation}
where $w = (w_s, w_m, w_l)$ are the butterfly weights.  The first two
zero constraints eliminate level and slope exposure; the third constraint
sets unit curvature sensitivity.

This system is solved via \texttt{numpy.linalg.solve} at each rolling
window.  When the loading matrix is near-singular (condition
number $> 10^{4}$), the date is skipped and no PnL is recorded for that
day.  In practice this never occurs in our sample; all 500
out-of-sample days produce valid weights.

\subsection{Daily PnL Computation}
The daily butterfly PnL is computed as:
\begin{equation}
\text{PnL}_t = w_{t-1}^{\!\top} \Delta y_t \times 10,000
\label{eq:pnl}
\end{equation}
where $\Delta y_t$ is the vector of rate changes at the three butterfly
tenors on day $t$ and $w_{t-1}$ are the weights estimated from the
previous day's rolling PCA.  The factor of $10{,}000$ converts decimal
rate changes to basis points.  Importantly, we use yesterday's weights
with today's rate changes to avoid look-ahead bias.

\subsection{Weight Dynamics}
Figure~\ref{fig:weights} illustrates the evolution of butterfly weights
for the spot 1Y--30Y strategy.  All three weights shift substantially
during regime transitions.  The belly weight (7Y) is not obviously more
stable than the wings: because the 7Y tenor carries the largest absolute
curvature loading (the trough of the U-shape), it is the most sensitive
weight to curvature PC rotations, while the wing weights are more
anchored by the stable level and slope constraints.

\begin{figure}[H]
\centering
\includegraphics[width=0.92\textwidth]{weight_stability_spot_1y-30y.png}
\caption{Butterfly weight dynamics for spot 1Y--30Y (3s/7s/15s).
(a)~Rolling weights over time.  (b)~Daily turnover
($\sum|\Delta w_i|$).  Turnover spikes coincide with regime transitions
and PC role reassignments.}
\label{fig:weights}
\end{figure}
The late-2024 turnover spike does not coincide cleanly with the
August~2024 rate cut.  That cut was decided by a narrow 5--4 MPC vote,
followed by an 8--1 hawkish hold in September.  The resulting
uncertainty about the pace of easing---rather than the cut itself---is
what shifted the realized covariance over subsequent months.

The April--May~2025 spike reflects compounding external shocks: the UK
Spring Budget (March~2025), the US tariff announcement (April~2025),
and UK local elections (May~2025), all of which altered the term
structure of uncertainty.

In both cases the proximate mechanism is eigenvector hypersensitivity:
when curvature and slope eigenvalues are nearly equal in the 252-day
window, small covariance perturbations produce large eigenvector
rotations and hence large weight jumps (Nadler,~2008).  This degeneracy
is most acute precisely when the weight instability occurs: the
full-sample $\lambda_2/\lambda_3$ ratio falls from 9.68 in the hold
regime to 1.34 in cutting, as simultaneous short-rate path and
term-premium uncertainty drive both eigenvalues up together, collapsing
the spectral gap.
% =====================================================================
\section{Trading Strategies}
\label{sec:strategies}
\subsection{Static Butterfly (Always-On)}
The static butterfly strategy is the simplest approach and holds the curvature portfolio continuously, re-estimating PCA loadings and recomputing butterfly weights each day using a rolling window. The portfolio is rebalanced daily to maintain zero exposure to level and slope and unit exposure to curvature, with P\&L given by the weighted change in rates at the three selected tenors. No entry or exit signal is applied, so the position remains continuously invested and simply adapts to the evolving curve structure.

\subsection{Z-Score Mean Reversion (Signal-Based)}
The signal strategy applies z-score mean reversion to the butterfly
spread level:
\begin{itemize}
  \item \textbf{Entry}: when the rolling 63-day z-score of the
        butterfly spread exceeds $\pm 1.5$ (spread is rich or cheap
        relative to recent history).
  \item \textbf{Exit}: when $|z| < 0.5$ (spread has reverted) or
        risk limits are hit.
  \item \textbf{Risk limits}: 50\,bps stop-loss and 30\,bps
        take-profit per trade.  The asymmetry is intentional: Leung \&
        Li~(2015) show that in an Ornstein--Uhlenbeck mean-reversion
        framework, a wider stop-loss implies a lower optimal take-profit,
        since the mean-reversion edge strengthens as the spread
        dislocates further, making premature exits costly.  The tight
        take-profit captures the reversion to fair value without
        overshoot risk.
\end{itemize}
The z-score exit is the primary mechanism; the bps limits are backstops
for cases where reversion fails.  The strategy is long curvature when
$z < -1.5$ and short when $z > +1.5$, trading 17--52 times over the
500-day out-of-sample period.  The 63-day window is a standard
quarterly convention (252/4), chosen to balance signal responsiveness
against noise.

\subsection{Literature-Inspired Variants}
We experiment with three variants that layer a systematic tilt on top of the
baseline PCA curvature butterfly.  Each is motivated by a well-known idea in
the fixed-income and systematic-trading literature, adapted here to the
specific setting of a PC3-weighted swap butterfly.  The goal is exploratory:
we test whether standard risk-management, timing, and signal-combination
techniques meaningfully alter the risk-adjusted performance of the base
strategy.

\subsubsection{Volatility-Scaled Butterfly}
Inverse-volatility position sizing is a standard risk-normalisation technique
in systematic trading (see, e.g., Moskowitz, Ooi, and Pedersen~(2012)) and
has been applied extensively to fixed-income carry and curve strategies.
In the context of PCA-based swap butterflies, Hariparsad and
Mar\'{e}~(2023) document substantial regime-dependent variation in butterfly
PnL volatility across hiking, cutting, and hold periods in the South African
swap market. This pattern is also observed in our own backtest. Motivated by
this, we scale the daily position size inversely proportional to trailing
63-day realised PnL volatility:
\begin{equation}
w_t^{\text{vol}} = w_t \cdot \frac{\sigma_{\text{target}}}{\hat\sigma_{t}},
\end{equation}
where $\sigma_{\text{target}} = 5$\,bps/day.  The rationale is specific to a
curvature butterfly: PC3 loadings amplify sensitivity to sharp
non-parallel moves during policy transitions, so unscaled position sizing
tends to produce outsized drawdowns in cutting cycles and under-utilised
capital in calm regimes.  This variant tests whether risk normalisation
alone, without any change to the entry signal, improves Sharpe.

\subsubsection{Curvature Momentum}
Time-series momentum in yield-curve factor returns has been documented
across markets.  Moskowitz, Ooi, and Pedersen~(2012) establish the
general time-series momentum effect across asset classes, including
fixed income, while subsequent practitioner work (e.g., Newfound
Research, 2019) finds evidence of positive autocorrelation in level,
slope, and curvature factor portfolios constructed from Treasury
futures.  The mechanism is plausible: monetary-policy cycles unfold
over months, segmented demand creates persistence at specific curve
points, and market expectations can adjust only gradually.

Given that our baseline strategy is a pure PC3 exposure, its PnL
provides a direct test of curvature momentum.  We implement a
21-day/5-day overlay: the position follows the sign of trailing 21-day
cumulative PnL and is rebalanced every 5~days.  If curvature shocks
persist over the lookback horizon, this variant should capture
trend-following behaviour; if curvature instead mean-reverts quickly,
it should underperform the baseline.

\subsubsection{Carry + Mean-Reversion Overlay}
Combining complementary signals to filter entries is a standard approach
in systematic trading.  Koijen, Moskowitz, Pedersen, and Vrugt~(2018)
identify carry as a robust cross-asset return predictor, while studies
such as Dreher, Gr\"{a}b, and Kostka~(2020) highlight the informational
content of curvature-related signals.  In the context of a PCA butterfly,
mean reversion captures short-horizon dislocations in the spread, whereas
the sign of trailing PnL acts as a proxy for carry, indicating whether the
position is supported by roll-down dynamics.

We combine a z-score mean-reversion signal (using a looser $|z|>1.0$
threshold) with a carry proxy given by the sign of trailing 21-day PnL:
\begin{itemize}
  \item Full position when both signals agree.
  \item Half position when only the MR signal fires.
  \item Flat when signals disagree, only carry fires, or neither fires.
\end{itemize}
This asymmetric scaling treats mean reversion as the primary signal and
carry as a confirming filter.  If carry adds incremental information,
the overlay should reduce losing trades without proportionately reducing
gains.
\section{Results}
\label{sec:results}

\subsection{Strategy Performance Overview}
Table~\ref{tab:summary} reports performance metrics across all three
curve configurations, including the static butterfly, z-score
mean-reversion signal, and three strategy variants.

\begin{table}[H]
\centering
\caption{Strategy Performance Summary.  Sharpe ratios are annualised
($\sqrt{252}$).  PnL and volatility are in basis points.  Hit rate is
the fraction of positive PnL days (lower for signal strategies that
are frequently flat).  Skewness is computed over all calendar trading
days; flat days contribute zero PnL rather than missing values, so the
many zeros in MR strategies reduce the standard deviation and
mechanically inflate the skewness statistic relative to an
always-in-market strategy.}
\label{tab:summary}
\small
\begin{tabular}{llcccccc}
\toprule
Curve & Strategy & Sharpe & PnL (bps) & Vol (bps) & MaxDD & Hit\% & Skew \\
\midrule
\multirow{5}{*}{spot 1Y--30Y}
  & Static       & 0.58 & 214  & 187 & 180 & 53.8 & 0.51 \\
  & Z-Score MR   & 1.14 & 285  & 126 & 83  & 25.0 & 0.68 \\
  & Vol-Scaled   & 0.92 & 127  & 79  & 62  & 27.9 & 1.55 \\
  & Momentum     & 1.18 & 431  & 188 & 78  & 53.2 & $-$0.20 \\
  & Carry+MR     & \textbf{1.31} & 352  & 136 & \textbf{77}  & 31.1 & 0.77 \\
\midrule
\multirow{5}{*}{spot 1Y--10Y}
  & Static       & \textbf{0.84} & 222  & 134 & \textbf{46} & 52.8 & 12.17 \\
  & Z-Score MR   & 0.76 & 182  & 121 & 39  & 22.2 & 16.39 \\
  & Vol-Scaled   & $-$0.23 & $-$32  & 79  & 80  & 24.0 & $-$3.18 \\
  & Momentum     & $-$0.86 & $-$226 & 135 & 266 & 49.3 & $-$11.92 \\
  & Carry+MR     & 0.42 & 62   & 76  & 48  & 28.9 & 7.21 \\
\midrule
\multirow{5}{*}{adjacent forwards}
  & Static       & 0.34 & 199  & 296 & 112 & 49.2 & 1.03 \\
  & Z-Score MR   & \textbf{1.42} & \textbf{539} & 191 & 121 & 19.2 & 1.27 \\
  & Vol-Scaled   & 0.39 & 136  & 202 & \textbf{94}  & 17.6 & 18.77 \\
  & Momentum     & $-$0.42 & $-$244 & 299 & 409 & 46.6 & 1.33 \\
  & Carry+MR     & 1.41 & 465  & 166 & 161 & 25.7 & $-$1.58 \\
\bottomrule
\end{tabular}
\end{table}

Key observations from Table~\ref{tab:summary}:
\begin{itemize}
  \item The z-score MR signal improves the Sharpe ratio on two of three
        curves, with the largest improvement on the forward curve
        (0.34~$\to$~1.42).  On spot~1Y--10Y the static butterfly
        already performs strongly (0.84) and MR provides no uplift
        (0.76).
  \item Carry+MR is the strongest variant on spot~1Y--30Y (1.31) and
      adjacent~forwards (1.41), narrowly below z-score MR on the forward
        curve while achieving the tightest max drawdown on 1Y--30Y (77\,bps).  It
        underperforms on spot~1Y--10Y (0.42), where the static strategy
        is already well-calibrated.
  \item Momentum works on spot~1Y--30Y (Sharpe 1.18) but \emph{reverses}
      on spot~1Y--10Y ($-$0.86) and adjacent~forwards ($-$0.42).  The
        spot~1Y--10Y reversal should be interpreted with caution: the
        near-50\% hit rate combined with a skew of $-$11.92 points to
        catastrophic losses on a small number of days, most likely
        the September~2024 BoE rate cut catching the momentum strategy
        on the wrong side.  This is event risk, not structural absence
        of curvature trending.
  \item Vol-scaling, which sizes z-score MR positions in proportion to a
        rolling volatility target, decreases the Sharpe ratio on
        spot~1Y--30Y (0.92 versus 1.14 for unrestricted MR) but reduces
        drawdown on the forward curve (94\,bps vs.\ 121\,bps for MR) as expected. It
        produces negative returns on spot~1Y--10Y ($-$0.23), where the
        MR signal itself is weak.  This is consistent with vol-scaling
        acting as a position sizer rather than an independent return source.
  \item The spot~1Y--10Y static strategy exhibits extreme positive
        skewness (12.17) driven by a single large gain during the
        September~2024 BoE rate cut.  Skew comparisons between Static
        and MR strategies are not fully like-for-like, since flat days
        in MR strategies mechanically inflate the skewness statistic.
\end{itemize}

\subsection{Cumulative PnL and Drawdowns}
Figure~\ref{fig:variant_pnl} shows cumulative PnL paths for all
variants on the forward-rate curve, the configuration with the widest
variant-to-variant Sharpe dispersion ($-$0.42 to 1.42).  The carry+MR
overlay achieves a Sharpe of 1.41, narrowly below the z-score MR
strategy (1.42) while reducing annualised volatility from 191 to
166\,bps.  Both PnL (539~$\to$~465\,bps) and volatility fall roughly
proportionally, so the carry gate primarily acts as a position filter
that reduces bet size rather than selectively avoiding losses; the
risk-adjusted improvement over z-score MR alone is minimal on this
curve.  The static strategy accumulates PnL slowly but steadily.
Volatility-scaling (136\,bps cumulative, Sharpe 0.39) sits
below the static in cumulative PnL terms; it reduces drawdown
relative to MR alone (121~$\to$~94\,bps) while retaining most of the
absolute return, consistent with vol-scaling acting as a position sizer
on the MR signal rather than an independent return source.
Momentum produces deep and sustained drawdowns.

\begin{figure}[H]
\centering
\includegraphics[width=0.92\textwidth]{variant_cumulative_pnl_fwd.png}
\caption{Cumulative PnL of all strategy variants on the forward-rate
butterfly (3s/7s/15s forward tenors).  Regime shading highlights the
hiking (red), hold (blue), and cutting (green) periods.}
\label{fig:variant_pnl}
\end{figure}

\begin{figure}[H]
\centering
\includegraphics[width=0.92\textwidth]{variant_cumulative_pnl_spot_1y-30y.png}
\caption{Cumulative PnL of all strategy variants on the spot 1Y--30Y
curve (3s/7s/15s butterfly).  Momentum performs strongly here, unlike
on other curves.}
\label{fig:variant_pnl_spot}
\end{figure}

\begin{figure}[H]
\centering
\includegraphics[width=0.92\textwidth]{variant_cumulative_pnl_spot_1y-10y.png}
\caption{Cumulative PnL of all strategy variants on the spot 1Y--10Y
curve (2s/5s/10s butterfly).  The dominant feature is a single large
gain during the September~2024 BoE rate cut announcement, which drives
the extreme positive skewness (12.17) of the static strategy.  Results
on this curve should be interpreted with caution.}
\label{fig:variant_pnl_spot10}
\end{figure}

Figure~\ref{fig:drawdown} presents the drawdown (underwater) chart
for the static strategies across all three curves.  The spot~1Y--30Y
strategy experiences the deepest drawdown (180\,bps) during mid-2024,
coinciding with the hold-to-cut regime transition.  The spot~1Y--10Y
strategy has remarkably shallow drawdowns (46\,bps) but with extreme
positive skewness driven by a single large gain during the
September~2024 rate cut.

\begin{figure}[H]
\centering
\includegraphics[width=0.92\textwidth]{drawdown_comparison.png}
\caption{Underwater chart showing drawdowns for the static butterfly
across all three curve configurations.  Regime shading as before.}
\label{fig:drawdown}
\end{figure}

\subsection{Factor Correlation Analysis}
\label{sec:factor_corr}
A well-constructed butterfly should be uncorrelated with level and slope
factors and strongly correlated with the curvature factor.
Table~\ref{tab:factor} reports the regression results from projecting
daily butterfly PnL onto the three PC factor scores.

\begin{table}[H]
\centering
\caption{Factor correlation analysis.  Regression:
$\text{PnL}_t = \alpha + \beta_1 \cdot \text{Level}_t
+ \beta_2 \cdot \text{Slope}_t + \beta_3 \cdot \text{Curvature}_t
+ \varepsilon_t$.}
\label{tab:factor}
\small
\begin{tabular}{lcccccc}
\toprule
Curve & $\rho_{\text{level}}$ & $\rho_{\text{slope}}$
  & $\rho_{\text{curv}}$ & $R^2$ & $\beta_{\text{curv}}$ \\
\midrule
spot 1Y--30Y & $-$0.065 & $+$0.003 & $+$0.209 & 0.057 & 0.855 \\
spot 1Y--10Y & $+$0.188 & $+$0.143 & $+$0.401 & 0.165 & 0.878 \\
adjacent forwards & $-$0.499 & $+$0.505 & $+$0.902 & \textbf{0.830} & 1.275 \\
\bottomrule
\end{tabular}
\end{table}

The forward-rate butterfly achieves an $R^{2}$ of 0.83 with the
three-factor model, compared to 0.06 for spot~1Y--30Y.  The curvature
beta for the forward-rate butterfly (1.275) is close to but exceeds
unity, suggesting mild \emph{over}-exposure to curvature relative to
the theoretical $\beta=1$ target.

The spot~1Y--30Y butterfly has near-zero level and slope correlations
($-$0.065 and $+$0.003), indicating effective hedging.  The spot~1Y--10Y
and forward-rate strategies show non-negligible cross-factor
correlations, which we discuss further in Section~\ref{sec:residual}.

Figure~\ref{fig:factor_corr} shows scatter plots of PnL against each
factor score for the forward-rate butterfly.  The curvature panel
exhibits a strong linear relationship, while level and slope show
moderate correlations that partially cancel in the regression.

\begin{figure}[H]
\centering
\includegraphics[width=0.92\textwidth]{factor_correlation_fwd.png}
\caption{Factor correlation scatter plots for the forward-rate
butterfly.  Strong curvature correlation ($\rho=0.90$) confirms
effective factor isolation.}
\label{fig:factor_corr}
\end{figure}

Figures~\ref{fig:factor_corr_spot30} and~\ref{fig:factor_corr_spot10} show the
equivalent scatter plots for the two spot-rate strategies.  The contrast is
striking: both spot curves produce diffuse curvature panels ($\rho \leq 0.40$)
compared to the tight linear relationship on the forward curve, directly
illustrating why forward-rate PCA achieves superior factor isolation.

\begin{figure}[H]
\centering
\includegraphics[width=0.92\textwidth]{factor_correlation_spot_1y-30y.png}
\caption{Factor correlation scatter plots for the spot 1Y--30Y butterfly.
Curvature correlation ($\rho=0.21$, $R^{2}=0.06$) is weak, confirming that
spot rates embed overlapping forward exposures that blur factor separation.}
\label{fig:factor_corr_spot30}
\end{figure}

\begin{figure}[H]
\centering
\includegraphics[width=0.92\textwidth]{factor_correlation_spot_1y-10y.png}
\caption{Factor correlation scatter plots for the spot 1Y--10Y butterfly.
Higher curvature correlation ($\rho=0.40$, $R^{2}=0.17$) than the full spot
curve but still far below the forward-rate result.}
\label{fig:factor_corr_spot10}
\end{figure}

\subsubsection{PC Role Tracking}
Figure~\ref{fig:pc_role} provides a compact view of which ordinal PC
is assigned to the curvature role across all three strategies.  The
spot~1Y--10Y strategy has the most evenly split assignment (curvature
is PC2 in $\approx$55\% of windows, PC3 in $\approx$45\%), making it
the least stable configuration.  The spot~1Y--30Y and forward-rate
strategies are somewhat more consistent ($\approx$65\% in one role),
though swapping still occurs in roughly a third of windows in both.
This confirms that ordinal PC labelling is unreliable across all
configurations and structural identification is essential.

\begin{figure}[H]
\centering
\includegraphics[width=0.92\textwidth]{pc_role_tracking.png}
\caption{PC role tracking: which ordinal PC (1, 2, or 3) is assigned
to the curvature role over time for each curve configuration.  Regime
shading as before.}
\label{fig:pc_role}
\end{figure}

% =====================================================================
\section{Discussion}
\label{sec:discussion}

\subsection{Residual Level and Slope Exposure}
\label{sec:residual}
The butterfly is constructed to have \emph{exactly} zero exposure to
the level and slope PCs \emph{as estimated in the rolling window}.
The non-zero factor correlations observed out of sample arise from
three sources:

\begin{enumerate}
  \item \textbf{Factor structure shifts:} When the BOE transitions
        between regimes, the covariance structure changes.  Weights
        calibrated using hold-regime data do not perfectly hedge the
        cutting-regime factor structure.
  \item \textbf{PC role swapping:} When the curvature assignment
        switches from PC3 to PC2 (or vice versa) between consecutive
        windows, the previous day's weights target the wrong
        eigenvector for one day, leaking level/slope exposure.
  \item \textbf{Instrument limitation:} A 3-tenor butterfly has exactly
        3 degrees of freedom.  Two are used to zero out level and slope;
        the third sets unit curvature.  There is no remaining freedom to
        hedge higher-order terms (e.g.\ a 4th PC explaining the residual
        0.1\% of variance).  With 9 maturities but only 3 instruments,
        the hedge is approximate even within the estimation window.
\end{enumerate}

The forward-rate implementation shows higher individual correlations
with level ($-$0.50) and slope ($+$0.51) than the spot~1Y--30Y strategy,
yet has a much higher curvature $R^{2}$ (0.83 vs.\ 0.06).  This
apparent paradox arises because forward-rate PCA factors are less
orthogonal to the spot-rate factors used in the regression.  The
forward-rate butterfly captures curvature movements very efficiently,
but its PnL also loads on level and slope \emph{measured in the
forward-rate basis}, which partially correlates with spot-rate level
and slope factors.

\subsection{Why Forward Rates Dominate}
Forward rates are differential quantities:
$f(T_i) \approx S(T_i) + T_i S'(T_i)$, which amplifies local curvature.
PCA on forward-rate changes therefore naturally separates local
curvature from the integrated (level) component that dominates
spot-rate changes.  This yields:
\begin{itemize}
  \item \textbf{Cleaner factor separation:} 90\% curvature correlation
        for forwards vs.\ 21\% for spots.
  \item \textbf{Better signal quality:} z-score MR on forwards achieves
        Sharpe 1.42 (the best single strategy in this sample) because curvature
        dislocations are more precisely measured.
  \item \textbf{More efficient hedging:} the forward-rate curvature
        loading is more localised, requiring less extreme wing weights
        and producing lower turnover.
\end{itemize}

\subsection{Strategy Variant Analysis}
The three variants test different aspects of curvature factor dynamics.
A striking feature of the results is the degree of cross-configuration
instability: no variant is uniformly superior across the three curve
specifications, and several flip sign entirely.  This is consistent with
Cochrane and Piazzesi's~(2005) finding that the yield-curve return-forecasting
factor is ``poorly related'' to the standard level, slope, and curvature
movements, implying that strategy performance is highly sensitive to how the
curve is parameterised.

\textbf{Carry+MR overlay} performs best on spot~1Y--30Y (Sharpe 1.31) and
adjacent~forwards (1.41), consistent with Dreher et al.'s~(2020)
finding that carry and value signals are complementary in curvature
trading.  However, it underperforms on spot~1Y--10Y (0.42 vs.\ static
0.84), confirming that the enhancement is not robust across configurations.

\textbf{Volatility scaling} sizes z-score MR positions by a rolling
volatility target.  It improves drawdown on the forward curve
(121~$\to$~94\,bps) but reduces absolute PnL sharply, as high volatility
in the cutting regime---precisely when MR is most profitable---causes
position sizes to be cut.  It produces negative returns on spot~1Y--10Y
($-$0.23), where the underlying MR signal is itself weak.

\textbf{Momentum} shows strong curve dependence.  It performs well on
spot~1Y--30Y (Sharpe 1.18) but reverses on spot~1Y--10Y ($-$0.86)
and adjacent~forwards ($-$0.42).  The spot~1Y--10Y result appears
event-driven: extreme negative skew ($-$11.92) with near-50\% hit rate
suggests losses are concentrated in the September~2024 BoE cut.  The
long-end outperformance may reflect persistent LDI flows, but with one
cycle in the sample this remains speculative.

\subsection{Limitations}
Several limitations should be noted:
\begin{enumerate}
  \item \textbf{Transaction costs:} All PnL figures are gross of
        transaction costs.  The daily rebalancing inherent in rolling
        PCA would incur non-trivial bid--ask costs in practice,
        particularly for the wing positions.
  \item \textbf{Short sample and regime concentration:} The out-of-sample
        period (500 days, $\approx$2 years) spans a single monetary-policy
        transition.  Conventional fixed-income backtesting guidance suggests
        a minimum of 3--5 years across multiple complete cycles for reliable
        inference.  Critically, virtually all PnL across
        every strategy is generated in the cutting regime (post-August 2024).
        Ang et al.'s~(2011) analysis of term-structure dynamics under
        monetary policy regime shifts shows that conventional and
        unconventional policy transmissions differ materially, making
        single-regime results non-generalizable.  All Sharpe ratios and
        variant comparisons should be interpreted as conditional on this
        specific macro path rather than as general evidence of strategy edge.
  \item \textbf{No model risk:} We assume the PCA factor structure is
        the ``true'' data-generating process.  In reality, factors may
        be time-varying and non-linear.
  \item \textbf{Signal parameters:} The 63-day z-score window and
        entry/exit thresholds ($\pm1.5$, 0.5) were set by convention and
        not optimised.  Alternative lookbacks may perform materially
        differently across curves and regimes.
  \item \textbf{Skewness:} The spot~1Y--10Y strategies exhibit extreme
        skewness (12.17), driven by a single outlier event.  Sharpe
        ratios may overstate risk-adjusted performance when the PnL
        distribution is highly non-Gaussian.
\item \textbf{Lack of standalone trading signal:} The PCA-based butterfly
      isolates curvature risk but does not, by itself, provide a
      predictive trading signal.  Without additional inputs such as
      valuation, macro, or flow-based indicators, the strategy represents
      a passive exposure to curvature rather than a systematic source of
      excess returns.
\end{enumerate}


% =====================================================================
\section{Conclusion}
\label{sec:conclusion}
This paper constructs PCA-based curvature butterfly strategies on GBP
OIS rates and evaluates their performance across three curve
configurations, two base trading rules, and three literature-inspired
variants over a three-year sample spanning hiking, hold, and cutting
regimes.

Our principal findings are:

\begin{enumerate}
  \item \textbf{Forward-rate PCA produces materially cleaner curvature
        trades than spot-rate PCA.}  The forward-rate butterfly achieves
        an $R^{2}$ of 0.83 with the three-factor model, compared to 0.06
        for the full spot curve.  This is because forward rates amplify
        local curvature and reduce the multicollinearity inherent in
        overlapping zero-coupon rates.

  \item \textbf{PC ordering is unstable across monetary-policy regimes.}
        Curvature was identified as PC3 in only $\sim$65\% of rolling
        windows on the full spot curve; it was PC2 in the remainder.
        Structural identification by loading shape, rather than ordinal
        rank, is essential for robust trading.

  \item \textbf{Signal-based mean reversion improves on the static butterfly
        in this sample.}  The z-score MR strategy achieves Sharpe ratios of
        1.14--1.42 across curves, compared to 0.34--0.84 for the static
        position.  However, virtually all gains are concentrated in the
        post-August 2024 cutting regime; the strategies are largely flat
        during the hiking and hold periods.  This makes inference fragile:
        the improvement reflects curvature-mean-reversion being rewarded in
        one specific macro environment, not across multiple cycles.

  \item \textbf{No variant is universally superior; results are highly
        curve-configuration-dependent.}  Carry+MR performs best on
        spot~1Y--30Y (1.31) and adjacent~forwards (1.41) but drops to
        0.42 on spot~1Y--10Y.  Momentum reverses in sign across curves.
        Vol-scaling turns negative on spot~1Y--10Y.  This cross-configuration
        instability is consistent with Cochrane and Piazzesi's~(2005)
        finding that yield-curve return-forecasting factors are poorly
        related to the standard PCA movements, and with the broader
        literature showing that fixed-income strategy performance is
        regime- and parameterisation-sensitive.

  \item \textbf{Momentum is profitable on the long end but reverses
        elsewhere.}  The reversal on spot~1Y--10Y ($-$0.86 Sharpe) is
        characterised by near-50\% hit rate and extreme negative skew
        ($-$11.92), consistent with event-driven loss rather than
        structural absence of curvature trending.  A single adverse
        event (the September~2024 BoE cut) likely accounts for most of
        the loss.  The structural LDI-flow hypothesis for spot~1Y--30Y
        outperformance remains speculative given the one-cycle sample.

  \item \textbf{Residual factor exposure is a structural limitation.}
        Non-zero level/slope correlations arise from regime-induced
        factor structure shifts, PC role swapping, and the inherent
        limitation of hedging a 9-dimensional curve with a 3-tenor
        butterfly.

\item \textbf{The PCA butterfly is a construction, not a standalone strategy.}
      While the PCA-based butterfly provides a clean and dynamically
      adaptive way to isolate curvature risk, it does not constitute a
      complete trading strategy in isolation.  In practice, such
      constructions are used as building blocks within broader
      relative-value frameworks, where positions are driven by macro
      views, valuation signals, or flow-based indicators.  Our results
      show that performance improvements arise primarily from the
      introduction of signal overlays (e.g.\ mean reversion and carry),
      highlighting that curvature exposure alone is insufficient to
      generate consistent excess returns.
\end{enumerate}

The most robust finding is structural: forward-rate PCA provides materially
cleaner curvature isolation than spot-rate PCA, and structural PC
identification by loading shape is essential for robust construction.
The performance results---while encouraging in this sample---should be
treated as hypothesis-generating.  The 500-day out-of-sample window
spans a single regime transition, and the literature on monetary policy
regime shifts (Ang et al., 2011) cautions strongly against generalising
single-regime backtest results.  A longer sample spanning multiple
tightening and easing cycles would be required to establish whether the
observed Sharpe ratios reflect genuine and persistent edge.

Future work could extend the framework to multi-currency settings,
incorporate transaction-cost models, and explore alternative estimation
approaches (e.g.\ exponentially weighted covariance matrices) for more
responsive regime adaptation.


% =====================================================================
\section*{References}

\begin{enumerate}[label={[\arabic*]},leftmargin=*,itemsep=2pt]
  \item Credit Suisse (2012).\ ``PCA Unleashed.''\ Interest Rate Strategy
        Research.
  \item Koijen, R.S.J., Moskowitz, T.J., Pedersen, L.H., and Vrugt, E.B.~(2018).\
        ``Carry.''\ \textit{Journal of Financial Economics}, 127(2), 197--225.
  \item Dreher, F., Gr\"{a}b, J., and Kostka, T.~(2020).\ ``From Carry
        Trades to Curvy Trades.''\ \textit{The World Economy}, 43(3),
        758--780.
  \item Hariparsad, S.\ and Mar\'{e}, E.~(2023).\ ``Examining Swap Butterfly
        Risk Premia in South Africa.''\ \textit{Investment Analysts
        Journal}, 52(3), 196--212.
  \item Hariparsad, S.\ and Mar\'{e}, E.~(2024).\ ``Optimal Monetary and
        Fiscal Policies to Maximise Non-Parallel Risk Premia in Sovereign
        Bond Markets.''\ \textit{Journal of Risk and Financial
        Management}, 17(11), 510.
  \item Lammi, M.~(2022).\ ``Yield Curve Arbitrage in the USD Interest
        Rate Swap Market: Returns and Risks in 2012--2021.''\ Master's
        Thesis, Aalto University School of Business.
  \item Litterman, R.\ and Scheinkman, J.~(1991).\ ``Common Factors
        Affecting Bond Returns.''\ \textit{The Journal of Fixed Income},
        1(1), 54--61.
  \item Oprea, A.~(2022).\ ``The Use of Principal Component Analysis (PCA)
        in Building Yield Curve Scenarios and Identifying Relative-Value
        Trading Opportunities on the Romanian Government Bond Market.''\
        \textit{Journal of Risk and Financial Management}, 15(6), 247.
  \item Suimon, Y., Sakaji, H., Izumi, K., and
        Matsushima, H.~(2020).\ ``Autoencoder-Based Three-Factor Model
        for the Yield Curve of Japanese Government Bonds and a
        Trading Strategy.''\ \textit{Journal of Risk and Financial
        Management}, 13(4), 82.
  \item Clarus Financial Technology~(2023).\ ``PCA and Swaps: What the
        Data Tells Us.''\ Clarus FT Blog.
  \item Leung, T.\ and Li, X.~(2015).\ ``Optimal Mean Reversion Trading
        with Transaction Costs and Stop-Loss Exit.''\ \textit{International
        Journal of Theoretical and Applied Finance}, 18(3), 1550020.
  \item Lord, R.\ and Pelsser, A.~(2007).\ ``Level-Slope-Curvature: Fact
        or Artefact?''\ \textit{Applied Mathematical Finance}, 14(2),
        105--130.
  \item Cochrane, J.H.\ and Piazzesi, M.~(2005).\ ``Bond Risk Premia.''\
        \textit{American Economic Review}, 95(1), 138--160.
  \item Ang, A., Boivin, J., Dong, S., and Loo-Kung, R.~(2011).\
        ``Monetary Policy Shifts and the Term Structure.''\
        \textit{Review of Economic Studies}, 78(2), 429--457.
\end{enumerate}

\end{document}
